% TOP_PROBLEM: {"id": 30005168, "title": "Iitaka Subadditivity Conjecture for Kodaira Dimension", "category_id": 5, "category": {"id": 5, "name": "algebraic_geometry", "display_name": "Algebraic Geometry", "description": "Geometric objects defined by polynomial equations.", "slug": "algebraic-geometry", "order_index": 5, "created_at": "2026-07-31T15:26:25.671Z"}, "status": "partially_solved"}
% FIRST_POSTED: 2026-10-06
% ATTEMPT_STATUS: solved
% OPENAI_CLAIM: {"schema_version":1,"source_repo":"https://github.com/openai/math","source_commit":"adc7f1241b42e322a6451854ab7e4b4c146bf78a","release_date":"2026-10-06","match":"stronger","resolution":"proved","independently_reviewed":false,"second_pass":"agreed","families":[{"family":"033","title":"Iitaka subadditivity, variation, and logarithmic additivity","subject":"Algebraic and complex geometry","manuscripts":[{"title":"Logarithmic Kodaira dimension and whole-fiber variation","pdf_path":"preprints/Logarithmic-Kodaira-dimension-and-whole-fiber-variation-September-26-2026/paper.pdf","readme_path":"preprints/Logarithmic-Kodaira-dimension-and-whole-fiber-variation-September-26-2026/README.md","date":"2026-09-26","theorem_ref":"thm:projective-variation (sections/01-introduction.tex:70); thm:main (sections/01-introduction.tex:38)","statement_sources":[{"path":"preprints/Logarithmic-Kodaira-dimension-and-whole-fiber-variation-September-26-2026/build/sections/01-introduction.tex","line":70,"label":"thm:projective-variation"},{"path":"preprints/Logarithmic-Kodaira-dimension-and-whole-fiber-variation-September-26-2026/build/sections/01-introduction.tex","line":38,"label":"thm:main"}]},{"title":"Orbifold and logarithmic Iitaka subadditivity","pdf_path":"preprints/Orbifold-and-logarithmic-Iitaka-subadditivity-September-26-2026/paper.pdf","readme_path":"preprints/Orbifold-and-logarithmic-Iitaka-subadditivity-September-26-2026/README.md","date":"2026-09-26","theorem_ref":"main:orbifold (sections/introduction.tex:135)","statement_sources":[{"path":"preprints/Orbifold-and-logarithmic-Iitaka-subadditivity-September-26-2026/build/sections/introduction.tex","line":135,"label":"main:orbifold"}]},{"title":"Projective Hodge lines and ordinary Iitaka subadditivity","pdf_path":"preprints/Projective-Hodge-lines-and-ordinary-Iitaka-subadditivity-September-27-2026/paper.pdf","readme_path":"preprints/Projective-Hodge-lines-and-ordinary-Iitaka-subadditivity-September-27-2026/README.md","date":"2026-09-27","theorem_ref":"thm:ordinary (sections/01-introduction.tex:48)","statement_sources":[{"path":"preprints/Projective-Hodge-lines-and-ordinary-Iitaka-subadditivity-September-27-2026/build/sections/01-introduction.tex","line":48,"label":"thm:ordinary"}]},{"title":"The reverse logarithmic Kodaira inequality and additivity","pdf_path":"preprints/The-reverse-logarithmic-Kodaira-inequality-and-additivity-September-26-2026/paper.pdf","readme_path":"preprints/The-reverse-logarithmic-Kodaira-inequality-and-additivity-September-26-2026/README.md","date":"2026-09-26","theorem_ref":"cor:additivity (sections/introduction.tex:51)","statement_sources":[{"path":"preprints/The-reverse-logarithmic-Kodaira-inequality-and-additivity-September-26-2026/build/sections/introduction.tex","line":51,"label":"cor:additivity"}]}],"lean":{"doc_path":"lean/docs/033.md","comparators":["lean/ComparatorChallenges/LogKodairaFiberNegative.json","lean/ComparatorChallenges/LogKodairaFiberNegative.lean"],"declarations":[],"covers_main_theorem":false},"reasoning_trace":null},{"family":"034","title":"Log abundance for compact Kähler spaces under logarithmic Iitaka subadditivity","subject":"Algebraic and complex geometry","manuscripts":[{"title":"Schnell fiber spaces and good canonical models","pdf_path":"preprints/Schnell-fiber-spaces-and-good-canonical-models-September-24-2026/paper.pdf","readme_path":"preprints/Schnell-fiber-spaces-and-good-canonical-models-September-24-2026/README.md","date":"2026-09-24","theorem_ref":"thm:canonical-good-models (sections/03-good-model-interface.tex:10); thm:schnell (sections/01-introduction.tex:29)","statement_sources":[{"path":"preprints/Schnell-fiber-spaces-and-good-canonical-models-September-24-2026/build/sections/03-good-model-interface.tex","line":10,"label":"thm:canonical-good-models"},{"path":"preprints/Schnell-fiber-spaces-and-good-canonical-models-September-24-2026/build/sections/01-introduction.tex","line":29,"label":"thm:schnell"}]}],"lean":null,"reasoning_trace":null}]}
\documentclass[11pt]{article}
\usepackage[margin=25mm]{geometry}
\usepackage{fontspec}
\setmainfont{Latin Modern Roman}
\usepackage{amsmath,amssymb,mathtools}
\usepackage{unicode-math}
\setmathfont{Latin Modern Math}
\usepackage{enumitem}
\usepackage{xurl,hyperref,bookmark}
\hypersetup{colorlinks=true,urlcolor=blue}
\setcounter{secnumdepth}{0}
\setlength{\parindent}{0pt}
\setlength{\parskip}{5pt}
\setlength{\emergencystretch}{3em}
\begin{document}
\section*{Iitaka Subadditivity Conjecture for Kodaira Dimension}

\subsection{OpenAI's claim}
Manuscript 1 (family 033, 2026-09-26): The logarithmic Iitaka--Viehweg lower bound is claimed for projective fibrations over a smooth quasi-projective base of nonnegative logarithmic Kodaira dimension, using variation of the entire generic fiber; Manuscript 2 (family 033, 2026-09-26): For compact manifolds in Fujiki class $\mathcal C$ carrying rational normal-crossing boundaries, the orbifold Iitaka lower bound is claimed, implying its ordinary and logarithmic specializations. Manuscript 3 (family 033, 2026-09-27): The ordinary Kodaira dimension inequality is claimed for projective connected-fiber morphisms in characteristic zero, without restricting total dimension or assuming nonnegative Kodaira dimension upstairs; Manuscript 4 (family 033, 2026-09-26): For reduced normal-crossing projective pairs that are smooth on every boundary stratum over the open base, OpenAI claims logarithmic Kodaira additivity, including the cases with negative-infinite dimension. Manuscript 5 (family 034, 2026-09-24): The Campana--Peternell Kodaira inequality and the specified Schnell fiber-space conclusion are claimed from a pseudo-effective difference between a canonical multiple and the chosen effective or pulled-back divisor.

\subsection{Mathematical statement of the claimed result}
\paragraph{Mathematical claim 1.} Let $f:X\to Y$ be a connected-fiber surjective morphism of smooth connected projective complex varieties, with smooth geometric generic fiber $F$ satisfying $\kappa(F)=0$. If some ample Cartier divisor $H$ on $Y$ and integer $m_0>0$ have $m_0K_X-f^*H$ pseudo-effective, OpenAI claims
\[
\kappa(X)=\dim Y.
\]
The manuscript also invokes its companion's canonical good-model theorem: every smooth connected projective complex $T$ with pseudo-effective $K_T$ has a $K_T$-negative birational contraction to a normal projective rational-factorial klt $V$ with $K_V$ semiample, and on a common smooth resolution $p:W\to T$, $q:W\to V$,
\[
p^*K_T=q^*K_V+E,\qquad E\ge0,\quad q_*E=0.
\]
The good-model statement is explicitly credited to the companion theorem, not presented as an independent proof here.

\paragraph{Mathematical claim 2.} For a projective connected-fiber surjective morphism $f:U\to V$ with $U$ and $V$ smooth, connected and complex quasi-projective. Assume $\overline\kappa(V)\ge0$. With geometric generic fiber $F$, OpenAI claims
\[
\overline\kappa(U)\ge\kappa(F)+\max\{\overline\kappa(V),\operatorname{Var}(f)\}.
\]
Logarithmic Kodaira dimension is $\kappa(\overline U,K_{\overline U}+D)$ on a smooth reduced-normal-crossing compactification. Birational variation minimizes $\operatorname{trdeg}_{\mathbb C}k$ over algebraically closed fields $k$ permitting a birational definition of the entire generic fiber. The projective specialization asserts $\kappa(X)\ge\kappa(F)+\max\{\kappa(Y),\operatorname{Var}(f)\}$ for a connected-fiber morphism $f:X\to Y$ between smooth connected projective varieties with $\kappa(Y)\ge0$.

\paragraph{Mathematical claim 3.} Let $X$ be smooth compact connected in Fujiki class $\mathcal C$ (bimeromorphic to a compact Kähler space), with rational effective simple-normal-crossing boundary $\Delta$ of coefficients in $[0,1]$. For a connected-fiber surjective holomorphic $f:X\to Y$ onto normal compact irreducible $Y$, and a very general smooth fiber $F$, OpenAI claims
\[
\kappa(X,K_X+\Delta)\ge\kappa(F,K_F+\Delta|_F)+\kappa(f\mid\Delta).
\]
The last term is the birationally invariant neat orbifold-base Kodaira dimension: for each base prime, the infimum of the multiplicities of its dominating pullback components weighted by reciprocal $1-\Delta$ coefficients determines boundary coefficient $1-1/m$. It uses $\kappa(\mathrm{point})=0$ and $(-\infty)+a=-\infty$, including $a=-\infty$.

\paragraph{Mathematical claim 4.} Over any algebraically closed characteristic-zero field, for a surjective projective morphism $f:X\to Z$ whose fibers are connected and whose source and target are smooth connected projective varieties and geometric generic fiber $F$, OpenAI claims
\[
\kappa(X)\ge\kappa(F)+\kappa(Z).
\]
Kodaira dimension is the maximal image dimension of positive complete pluricanonical systems, with $-\infty$ for empty systems; $(-\infty)+a=-\infty$ and a point has dimension zero. This is the manuscript's ordinary Iitaka theorem. Its referenced auxiliary Hodge-line input assumes a rationally polarized pure variation whose highest nonzero filtration is a line and unipotent boundary monodromy: its Deligne--Schmid extension is nef and is big precisely when the highest-line period map is generically immersive. That technical input is not the main subadditivity statement.

\paragraph{Mathematical claim 5.} Take a connected-fiber surjective morphism $f:X\to Y$ of smooth connected projective complex varieties, reduced simple-normal-crossing divisors $E,D$ with $\operatorname{Supp}(f^*D)\subseteq\operatorname{Supp}E$, and $V=Y\setminus D$. Require $X$ and each nonempty boundary intersection stratum of $E$, restricted over $V$, to be smooth over $V$. For a very general fiber $F$ and $E_F=E|_F$, OpenAI claims
\[
\kappa(X,K_X+E)=\kappa(Y,K_Y+D)+\kappa(F,K_F+E_F).
\]
The independent main theorem proves the upper inequality; equality uses the companion logarithmic subadditivity result. The zero divisors and $-\infty$ Kodaira dimensions are included. Smoothness is required over the original open base, with no assertion of smoothness above $D$.

\subsection{Outline of the argument}
Manuscript 1 (family 033, 2026-09-26): Canonical root covers isolate discriminant and moduli terms without losing pluriform orders. Flatness of the relevant Hodge lines is converted into birational constancy by singular-weight extension estimates and a big adjoint model; marked-family positivity and finite-cover descent then account for whole-fiber variation. Manuscript 2 (family 033, 2026-09-26): A relative Iitaka reduction first replaces the original fibers by the part for which general-type addition applies. Root-cover comparisons retain the boundary multiplicities, and positivity of the resulting period line supplies the missing contribution from the base.

Manuscript 3 (family 033, 2026-09-27): A canonical root cover of a Kodaira-zero fiber gives a one-dimensional top Hodge space. The argument distinguishes the full period image from the rank seen by that line, establishes the required positivity on the compactification, and applies the canonical bundle formula to recover subadditivity. Manuscript 4 (family 033, 2026-09-26): Iterated logarithmic cotangent and Higgs maps turn a selected pluriform into base sections with controlled coefficients. Period-image positivity and a canonical bundle formula supply the reverse inequality; combining it with subadditivity yields equality under the stated stratum-smoothness assumptions.

Manuscript 5 (family 034, 2026-09-24): The total-space canonical model, with its good-model property, supplies the geometric reduction. Cartier intersection theory and proper pushforward compare the base with the canonical image without requiring either map to factor through the other; transfer of sections then yields the Kodaira and fiber formulas.

\subsection{Scope relative to this problem}
The ordinary Iitaka theorem asserts $\kappa(X)\ge\kappa(F)+\kappa(Y)$ for every projective connected-fiber morphism between smooth connected varieties in characteristic zero. This includes the entire recorded complex fiber-space assertion. The variation result additionally replaces the base term by $\max\{\kappa(Y),\operatorname{Var}(f)\}$ when the base has nonnegative Kodaira dimension; a negative base term already makes the ordinary inequality automatic. The orbifold companion includes zero boundaries and extends to its specified compact complex setting. The additivity theorem needs smoothness over its open base, and the Schnell companion imposes an extra pseudo-effectivity hypothesis; those restricted companions are unnecessary for the unrestricted ordinary claim.

\subsection{Formal verification}
For family 033, the scope page documents the following selections. Only the case of a negative-infinite logarithmic fiber dimension is selected; it forces negative-infinite dimension upstairs and vanishing of positive pluriforms. Finite-dimensional fibers and the negative-base case remain outside the linked statement. Manuscript 1 is not among the papers selected by that page, so the family link does not formalize its main conclusion. It does not establish formal coverage of the complete problem clause. Manuscript 2 is not among the papers selected by that page, so the family link does not formalize its main conclusion. It does not establish formal coverage of the complete problem clause. Manuscript 3 is not among the papers selected by that page, so the family link does not formalize its main conclusion. It does not establish formal coverage of the complete problem clause. This entry does not assert that the selected declarations formalize every part of manuscript 4's main conclusion. Comparator files: LogKodairaFiberNegative.lean. Declarations: OAI.ReverseLogKodaira.SmoothProjectiveVariety.veryGenerally\_fiber\_negative\_branch.

Family 034 has no Lean scope page in this snapshot; no formalization of these manuscript results is listed.

This site has not executed the Lean files or checked their dependencies. These descriptions report the scope of the linked declarations, rather than a new verification of the proofs.

\subsection{Status caveat}
OpenAI attributes these results to an internal model. Its release notice cautions that results lacking full formalization can contain mistakes. This site has not independently assessed the proofs; a solved label records the scope of the reported claim.

\subsection{Sources}
\begin{itemize}
\item OpenAI, \emph{Logarithmic Kodaira dimension and whole-fiber variation}, 2026-09-26; citation key \texttt{OAI:Logarithmic-Kodaira-dimension-and-whole-fiber-variation-September-26-2026}. \href{https://github.com/openai/math/blob/adc7f1241b42e322a6451854ab7e4b4c146bf78a/preprints/Logarithmic-Kodaira-dimension-and-whole-fiber-variation-September-26-2026/paper.pdf}{PDF (GitHub)}; \href{https://github.com/openai/math/raw/adc7f1241b42e322a6451854ab7e4b4c146bf78a/preprints/Logarithmic-Kodaira-dimension-and-whole-fiber-variation-September-26-2026/paper.pdf}{PDF (direct)}; \href{https://github.com/openai/math/blob/adc7f1241b42e322a6451854ab7e4b4c146bf78a/preprints/Logarithmic-Kodaira-dimension-and-whole-fiber-variation-September-26-2026/README.md}{Manuscript page and citation}; \href{https://github.com/openai/math/blob/main/preprints/Logarithmic-Kodaira-dimension-and-whole-fiber-variation-September-26-2026/paper.pdf}{Latest version}; Statement source: \href{https://github.com/openai/math/blob/adc7f1241b42e322a6451854ab7e4b4c146bf78a/preprints/Logarithmic-Kodaira-dimension-and-whole-fiber-variation-September-26-2026/build/sections/01-introduction.tex\#L70}{thm:projective-variation}; Statement source: \href{https://github.com/openai/math/blob/adc7f1241b42e322a6451854ab7e4b4c146bf78a/preprints/Logarithmic-Kodaira-dimension-and-whole-fiber-variation-September-26-2026/build/sections/01-introduction.tex\#L38}{thm:main}.
\item OpenAI, \emph{Orbifold and logarithmic Iitaka subadditivity}, 2026-09-26; citation key \texttt{OAI:Orbifold-and-logarithmic-Iitaka-subadditivity-September-26-2026}. \href{https://github.com/openai/math/blob/adc7f1241b42e322a6451854ab7e4b4c146bf78a/preprints/Orbifold-and-logarithmic-Iitaka-subadditivity-September-26-2026/paper.pdf}{PDF (GitHub)}; \href{https://github.com/openai/math/raw/adc7f1241b42e322a6451854ab7e4b4c146bf78a/preprints/Orbifold-and-logarithmic-Iitaka-subadditivity-September-26-2026/paper.pdf}{PDF (direct)}; \href{https://github.com/openai/math/blob/adc7f1241b42e322a6451854ab7e4b4c146bf78a/preprints/Orbifold-and-logarithmic-Iitaka-subadditivity-September-26-2026/README.md}{Manuscript page and citation}; \href{https://github.com/openai/math/blob/main/preprints/Orbifold-and-logarithmic-Iitaka-subadditivity-September-26-2026/paper.pdf}{Latest version}; Statement source: \href{https://github.com/openai/math/blob/adc7f1241b42e322a6451854ab7e4b4c146bf78a/preprints/Orbifold-and-logarithmic-Iitaka-subadditivity-September-26-2026/build/sections/introduction.tex\#L135}{main:orbifold}.
\item OpenAI, \emph{Projective Hodge lines and ordinary Iitaka subadditivity}, 2026-09-27; citation key \texttt{OAI:Projective-Hodge-lines-and-ordinary-Iitaka-subadditivity-September-27-2026}. \href{https://github.com/openai/math/blob/adc7f1241b42e322a6451854ab7e4b4c146bf78a/preprints/Projective-Hodge-lines-and-ordinary-Iitaka-subadditivity-September-27-2026/paper.pdf}{PDF (GitHub)}; \href{https://github.com/openai/math/raw/adc7f1241b42e322a6451854ab7e4b4c146bf78a/preprints/Projective-Hodge-lines-and-ordinary-Iitaka-subadditivity-September-27-2026/paper.pdf}{PDF (direct)}; \href{https://github.com/openai/math/blob/adc7f1241b42e322a6451854ab7e4b4c146bf78a/preprints/Projective-Hodge-lines-and-ordinary-Iitaka-subadditivity-September-27-2026/README.md}{Manuscript page and citation}; \href{https://github.com/openai/math/blob/main/preprints/Projective-Hodge-lines-and-ordinary-Iitaka-subadditivity-September-27-2026/paper.pdf}{Latest version}; Statement source: \href{https://github.com/openai/math/blob/adc7f1241b42e322a6451854ab7e4b4c146bf78a/preprints/Projective-Hodge-lines-and-ordinary-Iitaka-subadditivity-September-27-2026/build/sections/01-introduction.tex\#L48}{thm:ordinary}.
\item OpenAI, \emph{The reverse logarithmic Kodaira inequality and additivity}, 2026-09-26; citation key \texttt{OAI:The-reverse-logarithmic-Kodaira-inequality-and-additivity-September-26-2026}. \href{https://github.com/openai/math/blob/adc7f1241b42e322a6451854ab7e4b4c146bf78a/preprints/The-reverse-logarithmic-Kodaira-inequality-and-additivity-September-26-2026/paper.pdf}{PDF (GitHub)}; \href{https://github.com/openai/math/raw/adc7f1241b42e322a6451854ab7e4b4c146bf78a/preprints/The-reverse-logarithmic-Kodaira-inequality-and-additivity-September-26-2026/paper.pdf}{PDF (direct)}; \href{https://github.com/openai/math/blob/adc7f1241b42e322a6451854ab7e4b4c146bf78a/preprints/The-reverse-logarithmic-Kodaira-inequality-and-additivity-September-26-2026/README.md}{Manuscript page and citation}; \href{https://github.com/openai/math/blob/main/preprints/The-reverse-logarithmic-Kodaira-inequality-and-additivity-September-26-2026/paper.pdf}{Latest version}; Statement source: \href{https://github.com/openai/math/blob/adc7f1241b42e322a6451854ab7e4b4c146bf78a/preprints/The-reverse-logarithmic-Kodaira-inequality-and-additivity-September-26-2026/build/sections/introduction.tex\#L51}{cor:additivity}.
\item Family 033: \href{https://github.com/openai/math/blob/adc7f1241b42e322a6451854ab7e4b4c146bf78a/lean/docs/033.md}{Lean scope}.
\item \href{https://github.com/openai/math/blob/adc7f1241b42e322a6451854ab7e4b4c146bf78a/lean/ComparatorChallenges/LogKodairaFiberNegative.json}{Comparator LogKodairaFiberNegative.json}.
\item \href{https://github.com/openai/math/blob/adc7f1241b42e322a6451854ab7e4b4c146bf78a/lean/ComparatorChallenges/LogKodairaFiberNegative.lean}{Comparator LogKodairaFiberNegative.lean}.
\item OpenAI, \emph{Schnell fiber spaces and good canonical models}, 2026-09-24; citation key \texttt{OAI:Schnell-fiber-spaces-and-good-canonical-models-September-24-2026}. \href{https://github.com/openai/math/blob/adc7f1241b42e322a6451854ab7e4b4c146bf78a/preprints/Schnell-fiber-spaces-and-good-canonical-models-September-24-2026/paper.pdf}{PDF (GitHub)}; \href{https://github.com/openai/math/raw/adc7f1241b42e322a6451854ab7e4b4c146bf78a/preprints/Schnell-fiber-spaces-and-good-canonical-models-September-24-2026/paper.pdf}{PDF (direct)}; \href{https://github.com/openai/math/blob/adc7f1241b42e322a6451854ab7e4b4c146bf78a/preprints/Schnell-fiber-spaces-and-good-canonical-models-September-24-2026/README.md}{Manuscript page and citation}; \href{https://github.com/openai/math/blob/main/preprints/Schnell-fiber-spaces-and-good-canonical-models-September-24-2026/paper.pdf}{Latest version}; Statement source: \href{https://github.com/openai/math/blob/adc7f1241b42e322a6451854ab7e4b4c146bf78a/preprints/Schnell-fiber-spaces-and-good-canonical-models-September-24-2026/build/sections/03-good-model-interface.tex\#L10}{thm:canonical-good-models}; Statement source: \href{https://github.com/openai/math/blob/adc7f1241b42e322a6451854ab7e4b4c146bf78a/preprints/Schnell-fiber-spaces-and-good-canonical-models-September-24-2026/build/sections/01-introduction.tex\#L29}{thm:schnell}.
\item \href{https://github.com/openai/math/blob/adc7f1241b42e322a6451854ab7e4b4c146bf78a/CONTENTS.md}{OpenAI catalogue entry}.
\end{itemize}
\end{document}
